\newpage
\section{Artifact Appendix}

%%%%%%%%%%%%%%%%%%%%%%%%%%%%%%%%%%%%%%%%%%%%%%%%%%%%%%%%%%%%%%%%%%%%%
\subsection{Abstract}

RAGMark is a hardware benchmarking framework for advanced Retrieval-Augmented Generation (RAG) pipelines.

This appendix describes how to access our code artifact and install it on any supported hardware.

To reproduce the figures in our paper, we provide access to a remote execution environment that allows reviewers to run the benchmark and generate data for three figures from our paper: Case 1 (naive, Fig.~9), Case 2 (reranking, Fig.~4), and Case 3 (compression, Fig.~5). See ``Hardware dependencies'' below for how this evaluation machine's hardware differs from the one used in the paper.




\subsection{Artifact check-list (meta-information)}

{\small
\begin{itemize}
  \item {\bf Program: } Advanced dense-retrieval RAG benchmarking tool supporting a wide range of configurations (embed, FAISS search, optional rerank/compression, LLM generation), and both standard and iterative pipelines.
  \item {\bf Hardware: } NVIDIA GPU. Paper results: 2$\times$ A100-SXM4-40GB. Artifact evaluation machine: 2$\times$ V100-PCIE-16GB. CPU/single-GPU subsets also supported.
  \item {\bf Metrics: } RAG Accuracy: (EM, F1, ROUGE-L, retrieval recall); Hardware performance (latency, TTFT, throughput, GPU power/energy, utilization).
  \item {\bf Output: } Per-query CSV/JSONL (accuracy, timing, energy, traces).
  \item {\bf Experiments: } Three reproducible experiments corresponding to paper figures: (1) reranking latency/accuracy (Fig.~4), (2) compression-rate latency (Fig.~5), (3) TTFT-vs-ROUGE trade-off (Fig.~9).
  \item {\bf How much disk space required (approximately)?: } $\sim$50--70GB (Wikipedia corpus + one FAISS index + generator weights).
  \item {\bf Publicly available?: } Yes.
  \item {\bf Code licenses (if publicly available)?: } MIT.
  \item {\bf Data licenses (if publicly available)?: } Wikipedia: CC BY-SA 3.0/GFDL; QA datasets via FlashRAG under original licenses.
  \item {\bf Archived (provide DOI)?: } [FILL IN once uploaded to Zenodo/Figshare]
\end{itemize}
}


%%%%%%%%%%%%%%%%%%%%%%%%%%%%%%%%%%%%%%%%%%%%%%%%%%%%%%%%%%%%%%%%%%%%%
\subsection{Description}

\subsubsection{How to access}
\url{https://github.com/zferic/RAGMark}, archived at DOI: [FILL IN once uploaded to
Zenodo/Figshare].
[NOTE: repo currently private — must be made public before submission so reviewers can
access it directly. IISWC's AE process is single-blind (reviewers see author identity),
so no anonymization is required on our end; see the SSH note below for how reviewer
anonymity is handled instead.]

\subsubsection{Hardware dependencies}
NVIDIA GPU is recommended, with hardware sized to the configuration being evaluated
(e.g., larger models/datasets, terabyte-scale databases, or CPU retrieval from SSD may
need multi-GPU); the framework is otherwise hardware-agnostic (NVIDIA or CPU).
Multi-GPU is also needed for the GPU0/GPU1 placement study (Sections 4.7-4.8).

Paper results were measured on 2$\times$ NVIDIA A100-SXM4-40GB GPUs, which the
university cannot provide remote access to. We instead provide 2$\times$ NVIDIA
V100-PCIE-16GB GPUs — an older, lower-memory generation. Absolute numbers (latency,
energy, memory) will not match the paper here, though they are reproducible on
matching hardware. We believe this also showcases RAGMark's portability across
hardware, and the impact both RAG configuration and hardware platform have on results.

\subsubsection{Software dependencies}

In our execution environment, all software dependencies necessary for operation are fully satisfied. Additionally, for AE evaluation, we use docker, but it is not a requirement outside of this ae environment prepared. All other dependencies are listed in requirements.txt and in the README.

\subsubsection{Data sets}
Retrieval Corpus: Wikipedia 2018 dump as retrieval corpus ($\sim$9.2M passages).
Evaluation QA Data: NQ, TriviaQA, SQuAD, WebQuestions, PopQA, HotpotQA (via FlashRAG). To recreate these datasets from scratch, we recommend using the FlashRAG toolkit at \url{https://github.com/RUC-NLPIR/FlashRAG} or your own.

\subsubsection{Models}
Supports any decoder model hosted on HuggingFace Hub — a HF token with LLM license acceptance is required (we configure this automatically in the execution environment for artifact evaluation.)

%%%%%%%%%%%%%%%%%%%%%%%%%%%%%%%%%%%%%%%%%%%%%%%%%%%%%%%%%%%%%%%%%%%%%
\subsection{Installation}
To install RAGMark from scratch on a new system, follow these steps:
\begin{verbatim}
git clone https://github.com/zferic/RAGMark.git
cd RAGMark

python -m venv .venv
source .venv/bin/activate

pip install -r requirements.txt
conda install -c pytorch -c nvidia faiss-gpu=1.8.0   # or faiss-cpu

# edit config.py: HF_HOME, BASE_OUT,
# WIKI_INDEX_DIR, WIKI_CORPUS_2018, DATASET_DIR

# build/place a FAISS index (Flat/IVF/IVF-SQ) under
# WIKI_INDEX_DIR via FlashRAG, plus the corpus JSONL
# and FlashRAG eval JSONL files under DATASET_DIR
\end{verbatim}

To simplify the evaluation process, we have prepared a remote execution environment
reachable via SSH, requiring only a single script to be run. Login credentials for a
shared reviewer account will be provided via HotCRP during the AE period, so reviewers
are not required to register using their own name, email, or institutional account.

\begin{verbatim}
ssh <username>@129.10.225.5
cd RAGMark
bash ae_evaluation/ae_eval.sh
\end{verbatim}

%%%%%%%%%%%%%%%%%%%%%%%%%%%%%%%%%%%%%%%%%%%%%%%%%%%%%%%%%%%%%%%%%%%%%
\subsection{Experiment workflow}

For this evaluation, we provide \texttt{ae\_evaluation/ae\_eval.sh} to load the Docker environment and directly produce the results end-to-end.

The script runs two sweep scripts followed by our analysis notebook (executed headlessly), which renders the figures below directly to \texttt{ae\_evaluation/figures/}:

\begin{itemize}
  \item {\bf Exp. 1 -- Case 2, Reranking Analysis (Section 4.3, Fig.~4): } \texttt{rerank=True}, \texttt{top\_ks=[0,1,3,5,10]}, all 3 generators $\times$ 6 datasets.
  \item {\bf Exp. 2 -- Case 3, Compression Analysis (Section 4.4, Fig.~5): } \texttt{compress=True}, \texttt{compress\_method="llmlingua2"}, \texttt{top\_ks=[1,3,5,10]}, \texttt{compress\_rates=[0.2,0.4,0.6,0.8]}.
  \item {\bf Exp. 3 -- Case 1, Naive Pipeline Analysis / Trade-offs (Sections 4.2 \& 4.8, Fig.~9): } naive pipeline (\texttt{rerank=False}, \texttt{compress=False}), \texttt{top\_ks=[0,1,3,5,10]}, all 3 generators.
\end{itemize}

Please note that RAG configuration support depends on the hardware support. A particular RAG configuration (e.g. 10 retrieved documents or a large LLM) may require substantially more compute than is available. Therefore, not all results that we ran on the A100 can be reproduced in the artifact evaluation V100 execution environment. 

%%%%%%%%%%%%%%%%%%%%%%%%%%%%%%%%%%%%%%%%%%%%%%%%%%%%%%%%%%%%%%%%%%%%%
\subsection{Evaluation and expected results}
\texttt{ae\_evaluation/load\_evals\_executed.ipynb} loads all output files into
\texttt{df\_results}/\texttt{df\_perf}/\texttt{df\_trace}, then reproduces the figures below and
writes each one to \texttt{ae\_evaluation/figures/plots/*.png}, alongside the exact aggregated
data table it was built from in \texttt{ae\_evaluation/figures/data/*.csv}.

\begin{itemize}
  \item {\bf Case 1 -- Naive Pipeline Analysis (Sections 4.2 \& 4.8, Fig.~9): } \texttt{df\_results} (ROUGE-L) merged with \texttt{df\_perf} (TTFT, decode, memory, SM\%), grouped by model $\times$ top-$k$, naive pipeline only.
  \item {\bf Case 2 -- Reranking Analysis (Section 4.3, Fig.~4): } \texttt{df\_perf} stage latencies (embed, FAISS, TTFT, decode) plus rerank duration recovered from \texttt{df\_trace}, merged with \texttt{df\_results} (ROUGE-L), grouped by dataset $\times$ model $\times$ top-$k$.
  \item {\bf Case 3 -- Compression Analysis (Section 4.4, Fig.~5): } \texttt{df\_perf} stage latencies (embed, FAISS, TTFT, decode) plus compression duration recovered from \texttt{df\_trace}, filtered to the compression sweep and grouped by top-$k$ $\times$ compression rate $r$.
\end{itemize}



%%%%%%%%%%%%%%%%%%%%%%%%%%%%%%%%%%%%%%%%%%%%%%%%%%%%%%%%%%%%%%%%%%%%%
\subsection{Methodology}
\begin{itemize}
  \item \url{https://www.acm.org/publications/policies/artifact-review-and-badging-current}
  \item \url{https://cTuning.org/ae}
\end{itemize}
